\subsection{Clifford swap対象の中心はcubic scrollとなる}
\label{subsec:normal-order2-scroll}
\begin{theorem}[偶Veroneseの全中心とquaternion order]
\label{thm:normal-order2-scroll-center}
$\Delta_h(b,c)\ne0$ のClifford swap H3反標準環を $R$ とする。
\[
 \begin{gathered}
 u=x+y,\quad v=x-y,\quad w=z+\frac h2u,\\
 B_0=v^2,\quad C_0=w^2,\quad D_0=\frac{uw+wu}{2},\quad a_0=vw,\quad s_0=uv+vu
 \end{gathered}
\]
と置くと、これら五つの二次元は $R$ の中心に属する。
それらの生成するgraded部分環は
\[
 \begin{gathered}
 \mathcal C=k[B_0,C_0,D_0,a_0,s_0]/
 (B_0C_0-a_0^2,\ C_0s_0-2D_0a_0,\ a_0s_0-2D_0B_0),\\
 \deg(B_0,C_0,D_0,a_0,s_0)=2
 \end{gathered}
\]
であり、$R$ は $\mathcal C$ 上有限・一般階数8である。
偶Veronese $R^{(2)}$ の全中心は、次数を半分にした $\mathcal C$ に等しい。
その中心上の一般階数は4、generic algebraはquaternion division algebraである。
特に $R$ のPI次数は2であり、
\[
 \operatorname{Proj}Z(R)=\operatorname{Proj}\mathcal C\simeq\mathbb F_1
\]
はcubic scrollである。
$\operatorname{qgr}R$ はこの $\mathbb F_1$ 上の階数4 orderのcoherent加群圏に同値となる。
この記述で $\tau^2$ はscrollのhyperplane線型束 $\OO_{\mathbb F_1}(E+2F)$ によるtensorに対応する。
orderの極大性・全分岐因子まではここで結論しない。
\end{theorem}
\begin{proof}
$T=t^2,P=xy+yx,Z=z^2$ は、中央Cliffordからのcocycle構成で$\Gamma$次数0の中心元の像なので
包絡の段階から中心である。
二次関係より
\[
 A_0=u^2=hZ+cT+P,\qquad B_0=hZ+cT-P,\qquad [v,w]=0,
\]
\[
 C_0=(1+h^3/4)Z+(3h^2/4)P+(hb+h^2c/4)T,
 \quad D_0=hP+(h^2/2)Z+(b+hc/2)T
\]
も中心である。交代tensorの式を $u,v,w$ に直すと
$G_*=-[vw,u]/4$ となる。
従って $a_0=vw$ は三次商で $u,v,w,t$ 全てと可換する。
$u^2,v^2$ が中心なので $s_0=uv+vu$ は $u,v$ と可換し、
$[v,w]=0$、$uw+wu=2D_0$、$vw=a_0$ を使えば $w$ とも可換する。
$t$ とは偶次数によって可換する。
同じ式から表示した三つのscroll minorが零となる。

$\alpha_0=(1+h^3)/(2h)$、$\beta_0=(2-h^3)/(4h)$、$\kappa=hb-c/h$ と置くと
\[
 C_0=\alpha_0A_0+\beta_0B_0+\kappa T,\qquad
 D_0=(3h/4)A_0-(h/4)B_0+bT.
\]
$A_0$ を消去した後の $D_0$ の $T$ 係数は
\[
 -\frac{(h^3-2)b-3hc}{2(1+h^3)}\ne0
\]
である。従って $T,P,Z$ と $B_0,C_0,D_0$ は可逆な線型変換で結ばれる。
五つの中心元を係数環に取れば、$u,v,w,t$ の全二乗と全反交換子はこの係数環に属する。
四生成元Cliffordの標準書換えにより $R$ はその中心部分環上有限となる。

抽象的な三minor環は行列
$\left(\begin{smallmatrix}C_0&a_0&2D_0\\a_0&B_0&s_0\end{smallmatrix}\right)$
のcubic scroll環であり、次元3の整な正規環である。
その次数を1に直したHilbert級数は $(1+2q)/(1-q)^3$ である。
具体的には
$C_0=\lambda r^2,a_0=\lambda rs,B_0=\lambda s^2,2D_0=\mu r,s_0=\mu s$
という標準scroll parametrizationから素性を得る。
三minorのHilbert--Burch解消からCohen--Macaulay性を得て、
頂点以外で滑らかなscrollの錐なので正規性が従う。
この環から $R$ の中心部分環への全射に追加核があれば次元は2以下となるが、
$R$ はその環上有限でGK次元3なので矛盾する。
従って表示した環が正確に $\mathcal C$ である。

Hilbert級数
\[
 H_R(q)=\frac{1-q^3}{(1-q)^4},\qquad
 H_{\mathcal C}(q)=\frac{1+2q^2}{(1-q^2)^3}
\]
の最高極の比は8である。
$K=\operatorname{Frac}\mathcal C$ と置き $C_0$ を逆にすれば $v=a_0w/C_0$ であり、
$R\otimes_{\mathcal C}K$ は $u,w,t$ による三変数Clifford環になる。
その二次形式は $u^2=A_0,w^2=C_0,t^2=T,uw+wu=2D_0$、他の反交換子は零である。
$C_0,D_0,a_0/C_0$ は一般点で独立な有理座標なので
$T(A_0C_0-D_0^2)\ne0$ となる。実際 $T$ の $D_0$ 係数は単元であり、
$A_0$ は $B_0,C_0,D_0$ の一次式だから $A_0C_0-D_0^2$ の $D_0^2$ 係数は $-1$ である。
一般階数8との比較により、次元8のこのClifford環からの全射は同型である。
その偶部分は中心 $K$ のquaternion algebraであり、$R$ が整域なので斜体となる。
全Clifford環の中心は二次体拡大で、中心上の次数は2である。
その生成元は実際に中央三次 $\Omega=(uw-D_0)t$ として $R$ に属し、
$\Omega^2=T(D_0^2-A_0C_0)$ を満たす。

$R^{(2)}$ の中心元は一般点で $K$ に属し、$R^{(2)}$ が $\mathcal C$ 上有限なので整である。
$\mathcal C$ の正規性からその元は $\mathcal C$ に属する。
従って偶Veroneseの全中心を得る。
$Z(R)^{(2)}=\mathcal C$ でもあるから、Projはcubic scroll $\mathbb F_1$ に一致する。
中心有限な偶Veroneseのsheaf化とVeronese同値によりorder表示を得る。
$R$ のshift2は $R^{(2)}$ のshift1であり、中心の標準hyperplaneによるtensorなので最後の主張も従う。
\end{proof}

指定例 $h=b=c=1$ では独立に、中心二次の次元5、十五の対称積の階数12、
三つのminorとvolume交換子恒等式を有理数上で確認した。
一般階数8・偶部分階数4のHilbert級数比と線型変換係数も
\texttt{verify\_cubic\_normal\_order2\_scroll.py} が厳密に検算する。
この中心計算はcocycle前後の圏を同一視することなく、得られたH3自体の中心的幾何を記述する。
